\section{Theoretische Mechanik}
\begin{definition}[Bezugssystem]
	Ort nur in Bezug auf einen anderen Ort
\end{definition}
Begriffe: Ortkoordinaten, Zeitkoordinate, Massepunkt\\
Ziel: Berechnung der Bahnkurve
\paragraph{Newtons Axiome}
\begin{enumerate}
	\item In Inertialsystemen wird die kräftefreie Bewegung durch konstante
	      Geschwindigkeit beschrieben werden.
	\item $\vec{F}=\odv{\vec{p}}{t}$
	\item $F_{actio}0-F_{reactio}$
\end{enumerate}
\paragraph{Erhaltungssätze}
\begin{itemize}
	\item Impulserhaltung $\odv{\vec{p}}{t}=0$
	\item Drehimpulserhaltung $\odv{\vec{L}}{t}=\odv{}{t} (\vec{r}\times\vec{p})=0$
	\item Energieerhaltungssatz: für konservative Kräfte ist Summe der Energien erhalten.
\end{itemize}
\begin{align*}
	\vec{F}_{\text{kons}}=-\nabla U
\end{align*}